\documentclass[10pt,a4paper]{amsart}
\usepackage[margin=19mm]{geometry}
\usepackage{amsmath,amssymb,mathtools}
\usepackage[scr=rsfs,cal=euler]{mathalfa}
\usepackage{xfrac}
\usepackage[unicode,hidelinks,bookmarks=false]{hyperref}
\newcommand{\ie}{i.e.\ }
\newcommand{\cf}{cf.\ }
\newcommand{\resp}{respectively\ }
\newcommand{\Subsection}{\subsection}
\providecommand{\nobreakdash}{\nobreak\hbox{-}}
\setlength{\emergencystretch}{2em}
\multlinegap0pt
\usepackage[shortlabels]{enumitem}
\usepackage{xcolor}
\newtheorem{theorem}{Theorem}
\newtheorem{lemma}{Lemma}
\newtheorem{claim}{Claim}
\newtheorem{observation}{Observation}
\title{Connectivity keeping paths in digraphs}
\author{Hojin Chu, Boram Park, Homoon Ryu}
\date{Source: arXiv:2608.25240v1 (CC BY 4.0)}
\begin{document}
\maketitle

\noindent\emph{Source and licence.} Connectivity keeping paths in digraphs. Version arXiv:2608.25240v1 (CC BY 4.0), distributed by arXiv under the Creative Commons Attribution 4.0 International licence, http://creativecommons.org/licenses/by/4.0/, as recorded in the arXiv licence metadata for this exact version and shown on the abstract page https://arxiv.org/abs/2608.25240v1. Full licence notice: \texttt{licenses/arxiv-2608.25240-CC-BY-4.0.txt}.\par\smallskip

\noindent\textit{Editorial scope.} Counted unit: the article's theorem thm:main (source0.tex lines 87--90). The article's complete original proof of it is delivered with it (source0.tex lines 247--372, one \texttt{proof} environment of 126 lines, which the article places away from the statement). No mathematics is written, repaired or re-derived: the statement and the proof are the article's own lines, in the article's own notation, and no retained line is substituted or re-flowed.  Retained uncounted as the source-local context the proof needs: lines 191--194.  Nothing was omitted from any retained span, so the package carries no omission record.  No retained line contains a citation, so the wrapper carries no bibliography and no bibliography entry.  No retained line contains a figure environment, an image-inclusion command or drawing code, so no image is delivered and the package declares no asset dependency.  Editorial formatting only, in addition: an AMS \texttt{amsart} preamble that restates the article's own declarations and macros used by the retained lines, namely the environments \texttt{theorem}, \texttt{lemma}, \texttt{claim}, \texttt{observation} on separate counters, together with the standard packages those lines need.  The retained environments print in this document as Theorem 1, Lemma 1, Claim 1, Observation 1 in order of appearance, whereas the article prints the same items with the numbering of its own full document.  The complete native source of the article as distributed in the arXiv e-print for this exact version is bundled unchanged as \texttt{sources/208560/source0.tex}.
\begin{lemma}\label{lem:leaving-H}
Let $D$ be a strongly connected digraph with $\delta^0(D)\ge m+1$, and let $(P,H)$ be an $m$-max pair of $D$ such that $V(H)\ne V(D)\setminus V(P)$.
Then $D$ has an $H$-leaving or an $H$-entering dipath of length at least $m+2$.
\end{lemma}

\begin{theorem}\label{thm:main}
Let $m$ be a positive integer.
Every strongly connected digraph $D$ with $\delta^0(D)\ge m+1$ contains a dipath $P$ of order $m$ such that $D-V(P)$ is strongly connected.
\end{theorem}

\begin{proof}[Proof of Theorem~\ref{thm:main}]
Let $D$ be a strongly connected digraph with $\delta^0(D)\ge m+1$ and let $(P,H)$ be an $m$-max pair of $D$.
If $V(H)=V(D)\setminus V(P)$, then $D-V(P)$ is strongly connected, so we are done.
Suppose to the contrary that $V(H)\ne V(D)\setminus V(P)$.
By Lemma~\ref{lem:leaving-H}, $D$ contains an $H$-leaving or an $H$-entering dipath of length at least $m+2$.
Since the conclusion of the theorem is invariant under reversing all arcs, we may assume that $D$ contains an $H$-leaving dipath of length at least $m+2$.

Among all pairs $(Q,Z)$ such that $Q$ is an $H$-leaving dipath of order $m+2$ and $Z$ is a $(V(Q)\setminus V(H),V(H))$-dipath, we choose one for which $\ell(Z)$ is maximum, and subject to this, for which the initial vertex of $Z$ has the largest possible index on $Q$.
Indeed, since $D$ is strongly connected, there exists a $(V(Q)\setminus V(H),V(H))$-dipath.
Write $Q=q_0q_1\cdots q_{m+1}$, so that $V(Q)\cap V(H)=\{q_0\}$.
Let $q_s$ be the initial vertex of $Z$, and let $z^*\in V(H)$ be its terminal vertex.
By the definition of a $(V(Q)\setminus V(H),V(H))$-dipath,
\begin{equation}\label{eq:Z-meets}
V(Z)\cap V(Q)\subseteq\{q_s,q_0\}
\qquad\text{and}\qquad
V(Z)\cap V(H)=\{z^*\}.
\end{equation}
Let $W=V(H)\cup V(Z)\cup V(Q)$.
Since $q_1\notin V(H)$, let $R=q_1r_2\cdots r_t$ be a longest dipath with initial vertex $q_1$ in $D-(W\setminus\{q_1\})$.
The case $s=1$, in which $q_1\in V(Z)$, is allowed.
By the choice of $R$,
\begin{equation}\label{eq:R-meets}
V(R)\cap W=\{q_1\}.
\end{equation}

\begin{claim}\label{clm:containment}
$N_D^+(r_t)\subseteq V(R)\cup\{q_1,\ldots,q_{m+1}\}$.
\end{claim}

\begin{proof}
Suppose that $y\in N_D^+(r_t)$ for some $y\in V(H)\cup(V(Z)\setminus\{q_s\})$.
Let
\[
B=
\begin{cases}
V(H)\cup V(R),&\text{if }y\in V(H),\\
V(H)\cup V(R)\cup V(Z[y,z^*]),&\text{if }y\in V(Z)\setminus\{q_s\}.
\end{cases}
\]
Since $H$ is strongly connected, $D[B]$ is strongly connected.
By~\eqref{eq:R-meets}, we have $V(R)\cap V(H)=\emptyset$, and hence $|B|\ge|V(H)|+t>|V(H)|$.
Moreover, $B\cap V(Q)\subseteq\{q_0,q_1\}$ by~\eqref{eq:Z-meets} and~\eqref{eq:R-meets}.
Thus $q_2q_3\cdots q_{m+1}$ is a dipath of order $m$ whose vertex set is disjoint from $B$.
It follows that $D-\{q_2,q_3,\ldots,q_{m+1}\}$ has a strong component of order greater than $|V(H)|$, which contradicts the maximality of $|V(H)|$.
Therefore $N_D^+(r_t)\cap\bigl(V(H)\cup(V(Z)\setminus\{q_s\})\bigr)=\emptyset$.
By the maximality of $R$, $r_t$ has no out-neighbor outside $V(R)$ in $D-(W\setminus\{q_1\})$.
Hence, by~\eqref{eq:Z-meets}, we have $N_D^+(r_t)\subseteq V(R)\cup\{q_1,\ldots,q_{m+1}\}$.
\end{proof}

Let $W_0:=V(H)\cup V(Z)\cup V(Q[q_0,q_s])$ be a subset of $W$.
Since $q_1\in W_0\setminus V(H)$, by the orientations of $Q$ and $Z$, the following is immediate.
\begin{observation}\label{obv:W}
The digraph $D[W_0]$ is strongly connected and $|W_0|>|V(H)|$.
\end{observation}
We now use the strong subdigraph $D[W_0]$ of $D$ to bound $t$ and $|N_D^+(r_t)\cap V(Q)|$.

\begin{claim}\label{clm:outneighbors}
The following statements hold.
\begin{enumerate}[label=\textup{(\roman*)},leftmargin=*,itemsep=1pt]
\item $2\le t\le m$.
\item $\bigl|N_D^+(r_t)\cap V(Q)\bigr|\ge m+3-t$.
\end{enumerate}
\end{claim}

\begin{proof}
If $t=1$, then $r_t=q_1$ and $V(R)=\{q_1\}$, so Claim~\ref{clm:containment} implies $d_D^+(r_t)\le m$, which contradicts $\delta^+(D)\ge m+1$.
Thus $t\ge 2$.
To show $t \le m$, suppose to the contrary that $t\ge m+1$.
Then $R':=r_2r_3\cdots r_{m+1}$ is a dipath of order $m$ whose vertex set is disjoint from $W_0$ by~\eqref{eq:R-meets}.
By Observation~\ref{obv:W}, the digraph $D-V(R')$ has a strong component of order greater than $|V(H)|$, which contradicts the choice of the $m$-max pair $(P,H)$.
Therefore (i) holds.

By Claim~\ref{clm:containment}, every out-neighbor of $r_t$ outside $V(Q)$ belongs to $\{r_2,\ldots,r_{t-1}\}$.
Since $d_D^+(r_t)\ge m+1$, it follows that $\bigl|N_D^+(r_t)\cap V(Q)\bigr|\ge d_D^+(r_t)-(t-2)\ge m+3-t$.
\end{proof}

By Claim~\ref{clm:outneighbors}(i), we have $m\ge2$.

\smallskip
\begin{claim}\label{clm:rotation}
$(r_t,q_j)\notin A(D)$ for every $j$ satisfying either
\[
2\le j\le\min\{s,t\} \qquad \text{or} \qquad s+1\le j\le t+1.
\]
\end{claim}

\begin{proof}
Suppose that $(r_t,q_j)\in A(D)$ for some $j$ with $2\le j\le t+1$. Set $b:=j+m-t$.
Since $t\le m$ and $j\le t+1$, $j\le b\le m+1$.
Since $t\ge2$, the dipath $R[r_2,r_t]$ is well-defined and its vertices are disjoint from $V(Q)$ by~\eqref{eq:R-meets}.
Let $R_0=R[r_2,r_t]+(r_t,q_j)+Q[q_j,q_b]$.
Then $|V(R_0)|=(t-1)+(b-j+1)=m$ and $V(R_0)$ is disjoint from $V(H)$.

We first consider the case $s+1\le j\le t+1$.
All vertices $q_j,\ldots,q_b$ have indices greater than $s$.
It follows from~\eqref{eq:Z-meets} and~\eqref{eq:R-meets} that $V(R_0)\cap W_0=\emptyset$.
By Observation~\ref{obv:W}, the digraph $D-V(R_0)$ has a strong component of order greater than $|V(H)|$, which contradicts the maximality of $|V(H)|$.

Now we consider the case $2\le j\le\min\{s,t\}$.
Then $b\le m$ and $Q':=Q[q_0,q_1]+(q_1,r_2)+R_0$ is an $H$-leaving dipath of order $m+2$ satisfying
\begin{equation}\label{eq:Qprime-in-Q}
V(Q')\cap V(Q)=\{q_0,q_1\}\cup\{q_j,\ldots,q_b\}.
\end{equation}
Since $s\ge j\ge2$, we have $q_1\ne q_s$, and hence $V(R)\cap V(Z)=\emptyset$ by~\eqref{eq:Z-meets} and~\eqref{eq:R-meets}.

Suppose $b<s$.
Let $Z'=Q[q_b,q_s]+Z$.
The vertices $q_{b+1},\ldots,q_s$ are disjoint from $V(Q')$ by~\eqref{eq:Qprime-in-Q}, while $V(Z)\cap V(Q')\subseteq\{q_0\}\subseteq V(H)$ by~\eqref{eq:Z-meets}.
Since $V(R)\cap V(Z)=\emptyset$, $Z'$ is a $(V(Q')\setminus V(H),V(H))$-dipath.
However, $\ell(Z')=(s-b)+\ell(Z)>\ell(Z)$, which contradicts the maximality of $\ell(Z)$.

Suppose $b\ge s$.
Since $j\le s\le b$, the vertex $q_s$ lies on $Q'$.
Moreover,~\eqref{eq:Z-meets},~\eqref{eq:Qprime-in-Q}, and $V(R)\cap V(Z)=\emptyset$ imply that $V(Z)\cap(V(Q')\setminus V(H))=\{q_s\}$ and $V(Z)\cap V(H)=\{z^*\}$.
Thus $(Q',Z)$ belongs to the collection of pairs over which $(Q,Z)$ was chosen.
Relabel $Q'=q_0'q_1'\cdots q_{m+1}'$, and let $s'$ be the index such that $q_{s'}'=q_s$.
Since $j\le t$, $s'=(t+1)+(s-j)>s$.
This contradicts the secondary maximality condition in the choice of $(Q,Z)$.
\end{proof}

By Claim~\ref{clm:outneighbors}(i), we have $2\le t\le m$.
Every $j\in\{2,\ldots,t\}$ satisfies either $j\le\min\{s,t\}$ or $s+1\le j\le t+1$.
Thus Claim~\ref{clm:rotation} implies that none of the $t-1$ vertices $q_2,\ldots,q_t$ is an out-neighbor of $r_t$.
Then, by Claim~\ref{clm:containment}, $N_D^+(r_t)\cap V(Q)\subseteq \{q_1, q_{t+1}, \ldots, q_{m+1}\}$ and so $|N_D^+(r_t)\cap V(Q)|\le 1+(m+1-t)=m+2-t$, which contradicts Claim~\ref{clm:outneighbors}(ii).
Hence, we have shown that $V(H)=V(D)\setminus V(P)$ and $D-V(P)$ is strongly connected.
\end{proof}

\end{document}
